\documentclass{beamer}
\usepackage[utf8]{inputenc}

\usetheme{Madrid}
\usecolortheme{default}
\usepackage{amsmath,amssymb,amsfonts,amsthm}
\usepackage{txfonts}
\usepackage{tkz-euclide}
\usepackage{listings}
\usepackage{adjustbox}
\usepackage{array}
\usepackage{tabularx}
\usepackage{gvv}
\usepackage{lmodern}
\usepackage{circuitikz}
\usepackage{tikz}
\usepackage{graphicx}

\setbeamertemplate{page number in head/foot}[totalframenumber]

\usepackage{tcolorbox}
\tcbuselibrary{minted,breakable,xparse,skins}



\definecolor{bg}{gray}{0.95}
\DeclareTCBListing{mintedbox}{O{}m!O{}}{%
	breakable=true,
	listing engine=minted,
	listing only,
	minted language=#2,
	minted style=default,
	minted options={%
		linenos,
		gobble=0,
		breaklines=true,
		breakafter=,,
		fontsize=\small,
		numbersep=8pt,
		#1},
	boxsep=0pt,
	left skip=0pt,
	right skip=0pt,
	left=25pt,
	right=0pt,
	top=3pt,
	bottom=3pt,
	arc=5pt,
	leftrule=0pt,
	rightrule=0pt,
	bottomrule=2pt,
	toprule=2pt,
	colback=bg,
	colframe=orange!70,
	enhanced,
	overlay={%
		\begin{tcbclipinterior}
			\fill[orange!20!white] (frame.south west) rectangle ([xshift=20pt]frame.north west);
	\end{tcbclipinterior}},
	#3,
}
\lstset{
	language=C,
	basicstyle=\ttfamily\small,
	keywordstyle=\color{blue},
	stringstyle=\color{orange},
	commentstyle=\color{green!60!black},
	numbers=left,
	numberstyle=\tiny\color{gray},
	breaklines=true,
	showstringspaces=false,
}
%------------------------------------------------------------
%This block of code defines the information to appear in the
%Title page
\title %optional
{5.2.49}
%\subtitle{A short story}

\author % (optional)
{RAVULA SHASHANK REDDY - EE25BTECH11047}

 \begin{document}
	
	
	\frame{\titlepage}
	\begin{frame}{Question}
     Find all circles that touch the $x$-axis at a distance $3$ from the origin and have a $y$-axis intercept of length $2\sqrt{7}$ .\\
\end{frame}
\begin{frame}{Solution}
    \begin{align}
\vec{x}^\top \vec{x} + 2 \vec{u}^\top \vec{x} + f &= 0, \quad \vec{C} = -\vec{u}, \quad r^2 = \vec{C}^\top \vec{C} - f \\
\vec{e}_1 &= \myvec{1 \\ 0}, \quad \vec{e}_2 = \myvec{0 \\ 1}\\
\quad \vec{C}&=\myvec{\vec{e}_1^\top \vec{C} \\ \vec{e}_2^\top \vec{C}}
\end{align}

Substituting y-axis in circle equation
\begin{align}
(y \vec{e}_2)^\top (y \vec{e}_2) + 2 \vec{u}^\top (y \vec{e}_2) + f &= 0 
\\ \implies \quad 
y^2 - 2 (\vec{C}^\top \vec{e}_2) y + f &= 0 \label{eq:m}
\end{align}

Here the square of Intercept equals to determinant of \eqref{eq:m}
\begin{align}
\Delta &= [-2 (\vec{C}^\top \vec{e}_2)]^2 - 4 f = 4 \big[ (\vec{C}^\top \vec{e}_2)^2 - f \big]=(2\sqrt{7})^2
\end{align}
\begin{align}
(\vec{C}^\top \vec{e}_2)^2 - f &= 7 \label{eq:e}
\end{align}
\end{frame}
\begin{frame}{Solution}
\begin{align}
\vec{C}^\top \vec{C} - f&=r^2 \label{eq:r}\\
\eqref{eq:r}-\eqref{eq:e}\implies\vec{C}^\top \vec{C} - (\vec{C}^\top \vec{e}_2)^2 &= r^2 - 7 
\quad \implies \quad (\vec{e}_1^\top \vec{C})^2 = r^2 - 7 
\end{align}

Since Circle touches x-axis at a distance 3 from origin. 
\begin{align}
|\vec{e}_2^\top \vec{C}|&=r\\
|\vec{e}_1^\top \vec{C}| &= 3 \quad \implies r^2 = 16 \\
(\vec{e}_2^\top \vec{C})^2 - f &= 7, \quad \implies f=9
\end{align}
\begin{align}
\vec{C} = \myvec{\pm3 \\ \pm4}
\end{align}
\end{frame}
\begin{frame}{Solution}
\begin{align}
\vec{x}^\top \vec{x} + 2(-\vec{C})^\top \vec{x} + f = 0 
\end{align}
\begin{align}
\vec{x}^\top \vec{x} + 2 \myvec{-3 & -4} \vec{x} + 9 = 0, \quad
\vec{x}^\top \vec{x} + 2 \myvec{-3 & 4} \vec{x} + 9 = 0, \\
\vec{x}^\top \vec{x} + 2 \myvec{3 & -4} \vec{x} + 9 = 0, \quad
\vec{x}^\top \vec{x} + 2 \myvec{3 & 4} \vec{x} + 9 = 0 
\end{align}
\end{frame}
\begin{frame}[fragile]
\frametitle{C Code}
\begin{lstlisting}
    #include <stdio.h>
#include <math.h>

int main() {
    double a = 3.0;
    double L = 2.0 * sqrt(7.0);
    double r = sqrt(a*a + (L/2)*(L/2));
    double f = a*a + (L/2)*(L/2) - r*r;

    double centers[4][2] = {
        { a,  r},
        { a, -r},
        {-a,  r},
        {-a, -r}
    };
    \end{lstlisting}  
\end{frame}
\begin{frame}[fragile]
\frametitle{C Code}
\begin{lstlisting}
    for(int i=0; i<4; i++) {
        double h = centers[i][0];
        double k = centers[i][1];
        printf("Circle %d: Centre = (%.2f, %.2f), Radius = %.2f\n", i+1, h, k, r);
    }

    printf("Constant term f = %.2f\n", f);
    return 0;
}

\end{lstlisting}
\end{frame}
\begin{frame}[fragile]
\frametitle{Python Direct}
\begin{lstlisting}
    import math
import numpy as np
import matplotlib.pyplot as plt
from funcs import circ_gen, circ_param

a = 3.0
L = 2.0 * math.sqrt(7.0)
half = L / 2.0
r = math.sqrt(a*a + half*half)
f = a*a

centers = [
    np.array([[ a],[ r]], dtype=float),
    np.array([[ a],[-r]], dtype=float),
    np.array([[-a],[ r]], dtype=float),
    np.array([[-a],[-r]], dtype=float)
]
\end{lstlisting}
\end{frame}
\begin{frame}[fragile]
\frametitle{Python Direct}
\begin{lstlisting}
for i, C in enumerate(centers, 1):
    h = float(round(C[0,0],2))
    k = float(round(C[1,0],2))
    R = float(round(r,2))
    print(f"Circle {i}: Centre = ({h}, {k}), Radius = {R}")

fig, ax = plt.subplots(figsize=(7,7))
ax.set_aspect('equal', adjustable='box')
ax.grid(True, linestyle='--', alpha=0.5)

for i, C in enumerate(centers,1):
    circ = circ_gen(C, r)
    ax.plot(circ[0,:], circ[1,:], linewidth=1.8)
    ax.scatter(float(C[0,0]), float(C[1,0]), color='black', s=25)
    dx = -0.6 if C[0,0] >= 0 else 0.6
    dy = -0.4 if C[1,0] >= 0 else 0.4
    h = float(round(C[0,0],2))
    k = float(round(C[1,0],2))
    \end{lstlisting}
\end{frame}
\begin{frame}[fragile]
\frametitle{Python Direct}
\begin{lstlisting}
    ax.text(h + dx, k + dy, f"O{i}({h},{k})", fontsize=10, bbox=dict(boxstyle="round,pad=0.2", fc="white", alpha=0.8))

for s in (a, -a):
    ax.scatter(s, 0, color='red', s=50, marker='x')
    ax.text(s + 0.25, -0.45, f"({s},0)", color='red')

ax.axhline(0, color='gray', linewidth=1)
ax.axvline(0, color='gray', linewidth=1)
ax.set_xlim(-8,8)
ax.set_ylim(-8,8)
ax.set_xlabel('x')
ax.set_ylabel('y')
ax.set_title('Circles tangent to x-axis at (±3,0) with y-axis intercept 2√7')
plt.show()


\end{lstlisting}
\end{frame}
\begin{frame}[fragile]
\frametitle{Python Shared}
\begin{lstlisting}
    import ctypes
import numpy as np

lib = ctypes.CDLL("./circles.so")

lib.compute_circles.argtypes = [
    ctypes.c_double,
    ctypes.c_double,
    np.ctypeslib.ndpointer(dtype=np.float64, shape=(4,2)),
    ctypes.POINTER(ctypes.c_double),
    ctypes.POINTER(ctypes.c_double)
]

lib.compute_circles.restype = None
\end{lstlisting}
\end{frame}
\begin{frame}[fragile]
\frametitle{Python Shared}
\begin{lstlisting}
a = 3.0
L = 2*np.sqrt(7)
centers = np.zeros((4,2), dtype=np.float64)
r = ctypes.c_double()
f = ctypes.c_double()

lib.compute_circles(a, L, centers, ctypes.byref(r), ctypes.byref(f))

print("Results from shared library:\n")
for i in range(4):
    h, k = centers[i]
    print(f"Circle {i+1}: Centre = ({round(h,2)}, {round(k,2)}), Radius = {round(r.value,2)}")

print(f"\nConstant term f = {round(f.value,2)}")

\end{lstlisting}
\end{frame}
\begin{frame}{Plot}
    \begin{figure}
        \centering
        \includegraphics[width=0.5\linewidth]{figs/Fig_1.png}
        \caption{}
        \label{fig:placeholder}
    \end{figure}
\end{frame}
\end{document}
